\Needspace{9\baselineskip}\begin{center}\textit{A method of successively reducing the degree of equations for trigonometric functions.}\end{center}


\Needspace{6\baselineskip}\Article{363}
Just as, on the assumption \(n-1=ef\), the function \(X\) can be resolved into \(e\) factors of degree \(f\) as soon as the values of all the \(e\) sums of \(f\) terms are known (Article 348), so too, if we suppose \(Z=0\) to be an equation of degree \(n-1\) whose roots are the sines or any other trigonometric functions of the angles \(\frac Pn,\allowbreak \frac{2P}n\ldots\frac{(n-1)P}n\), the function \(Z\) can then be resolved into \(e\) factors of degree \(f\). The principal points of this procedure are as follows.

Let \(\Omega\) consist of the \(e\) periods of \(f\) terms \((f,1)=P,\allowbreak P'.P''\), etc.; let the period \(P\) consist of the roots \([1],\allowbreak [a],\allowbreak [b],\allowbreak [c]\), etc.; \(P'\) of \([a'],\allowbreak [b'],\allowbreak [c']\), etc.; \(P''\) of \([a''],\allowbreak [b''],\allowbreak [c'']\), etc., etc. Let the root \([1]\) correspond to the angle \(\omega\), and therefore the roots \([a],\allowbreak [b]\), etc., to the angles \(a\omega,\allowbreak b\omega\), etc., the roots \([a'],\allowbreak [b']\), etc., to the angles \(a'\omega,\allowbreak b'\omega\), etc., and the roots \([a''],\allowbreak [b'']\), etc., to the angles \(a''\omega,\allowbreak b''\omega\), etc., etc. It is easy to see that all these angles taken together agree with the angles \(\frac Pn,\allowbreak \frac{2P}n,\allowbreak \frac{3P}n\ldots\frac{(n-1)P}n\) as regards their trigonometric functions.\HistoricalNote{*)}{Two angles agree in this respect if their difference equals either the whole circumference or some multiple of it. We could call such angles \textit{congruent modulo the circumference}, if we wished to understand congruence in a somewhat broader sense.} If, therefore, the function in question is denoted by the symbol \(\varphi\) placed before the angle, and we set the product of the \(e\) factors \[x-\varphi\omega,\quad x-\varphi a\omega,\quad x-\varphi b\omega\quad\text{etc.}\] \(=Y\), the product of \(x-\varphi a'\omega,\allowbreak x-\varphi b'\omega\), etc., \(=Y'\), and the product of \(x-\varphi a''\omega,\allowbreak x-\varphi b''\omega\), etc., \(=Y''\), and so on, we necessarily have the product \(YY'Y''\ldots=Z\). It now remains to show that all the coefficients in the functions \(Y,\allowbreak Y',\allowbreak Y''\), etc., can be reduced to a form such as
\SourceLeaf{463}
\[A+B(f,1)+C(f,g)+D(f,gg)\ldots+L(f,g^{e-1})\] Once this is done, they must evidently all be regarded as known as soon as the values of all the sums of \(f\) terms are known. We will accomplish this as follows.

Just as \(\cos\omega=\tfrac12[1]+\tfrac12[1]^{n-1},\allowbreak \ \sin\omega=-\tfrac12i[1]+\tfrac12i[1]^{n-1}\), so too, by the preceding article, the remaining trigonometric functions of the angle \(\omega\) can be reduced to a form such as \(\mathfrak A+\mathfrak B[1]+\mathfrak C[1]^2+\mathfrak D[1]^3+\text{etc.}\). It is readily seen that the function of the angle \(k\omega\) then becomes \(=\mathfrak A+\mathfrak B[k]+\mathfrak C[k]^2+\mathfrak D[k]^3+\text{etc.}\), where \(k\) denotes any integer. Now each coefficient in \(Y\) is a symmetric polynomial in \(\varphi\omega,\allowbreak \varphi a\omega,\allowbreak \varphi b\omega\), etc. It is therefore clear that substituting the values of these quantities makes each coefficient a symmetric polynomial in \([1],\allowbreak [a],\allowbreak [b]\), etc.; hence, by Article 347, they reduce to the form \(A+B(f,1)+C(f,g)+\text{etc.}\). By exactly the same reasoning, all the coefficients in \(Y',\allowbreak Y''\), etc., can also be reduced to a similar form. \textit{Q. E. D.}

\Needspace{6\baselineskip}\Article{364}
We add some further observations on the problem of the preceding article.

I. Since each coefficient in \(Y'\) is a function of the roots contained in the period \(P'\) (which we may set \(=(f,a')\)) identical to the function of the roots in \(P\) given by the corresponding coefficient in \(Y\), it is evident from Article 347 that \(Y'\) can be derived from \(Y\) by substituting everywhere in \(Y\), for \((f,1),\allowbreak (f,g),\allowbreak (f,gg)\), etc., respectively, \((f,a'),\allowbreak (f,a'g),\allowbreak (f,a'gg)\), etc. Similarly, \(Y''\) is derived from \(Y\) by substituting everywhere in \(Y\), for \((f,1),\allowbreak (f,g),\allowbreak (f,gg)\), etc., respectively, \((f,a''),\allowbreak (f,a''g),\allowbreak (f,a''gg)\), etc., etc. Thus, as soon as the function \(Y\) has been expanded, the remaining functions \(Y',\allowbreak Y''\), etc., follow from it without difficulty.

II. If we suppose \[Y=x^f-\alpha x^{f-1}+\beta x^{f-2}-\text{etc.}\] the coefficients \(\alpha,\allowbreak \beta\), etc., will respectively be the sum of the roots of the equation \(Y=0\quad \text{that is}\) the quantities \(\varphi\omega,\allowbreak \varphi a\omega,\allowbreak \varphi b\omega\), etc., the sum of their products in pairs, and so on. But these coefficients will usually be obtained much more conveniently by a method similar to that given in Article 349: calculating the sum of the roots \(\varphi\omega,\allowbreak \varphi a\omega,\allowbreak \varphi b\omega\), etc., the sum of their squares,
\SourceLeaf{464}
cubes, etc., and deriving those coefficients from these by Newton's theorem. — Whenever \(\varphi\) denotes the tangent, secant, cotangent or cosecant, there are further shortcuts, which we pass over here without discussion.

III. The case where \(f\) is even deserves special consideration; each period \(P,\allowbreak P',\allowbreak P''\), etc., is then composed of \(\tfrac12f\) periods of two terms. Let \(P\) consist of \((2,1),\allowbreak (2,\mathfrak{a}),\allowbreak (2,\mathfrak{b}),\allowbreak (2,\mathfrak{c})\), etc. The numbers \(1,\allowbreak \mathfrak{a},\allowbreak \mathfrak{b},\allowbreak \mathfrak{c}\), etc., and \(n-1,\allowbreak n-\mathfrak{a}.n-\mathfrak{b},\allowbreak n-\mathfrak{c}\), etc., taken together will then agree with \(1,\allowbreak a,\allowbreak b,\allowbreak c\), etc., or at least (which here amounts to the same thing) be congruent to them modulo \(n\). But we have \(\varphi(n-1)\omega=\pm\varphi\omega,\allowbreak \ \varphi(n-\mathfrak{a})\omega=\pm\varphi \mathfrak{a}\omega\), etc., the upper signs applying whenever \(\varphi\) denotes the cosine or secant, and the lower signs when \(\varphi\) denotes the sine, tangent, cotangent or cosecant. It follows that in the first two cases the factors making up \(Y\) always occur in equal pairs, so that \(Y\) is a square, and specifically \(Y=yy\) if \(y\) is set equal to the product of \[x-\varphi\omega,\ x-\varphi \mathfrak{a}\omega,\ x-\varphi \mathfrak{b}\omega\text{ etc.}\]

Similarly, in the same cases the remaining functions \(Y',\allowbreak Y''\), etc., will be squares. Specifically, if we suppose that \(P'\) consists of \((2,\mathfrak{a}'),\allowbreak (2,\mathfrak{b}'),\allowbreak (2,\mathfrak{c}')\), etc., \(P''\) of \((2,\mathfrak{a}''),\allowbreak (2,\mathfrak{b}''),\allowbreak (2,\mathfrak{c}'')\), etc., etc., and that the product of \(x-\varphi \mathfrak{a}'\omega,\allowbreak x-\varphi \mathfrak{b}'\omega,\allowbreak x-\varphi \mathfrak{c}'\omega\), etc., is \(=y'\), the product of \(x-\varphi \mathfrak{a}''\omega,\allowbreak x-\varphi \mathfrak{b}''\omega\), etc., \(=y''\), and so on, we will have \(Y'=y'y',\allowbreak \ Y''=y''y''\), etc. The function \(Z\) will also be a square (compare Article 337 above), and its square root will equal the product of \(y,\allowbreak y',\allowbreak y''\), etc. Further, it is easy to see that \(y',\allowbreak y''\), etc., are derived from \(y\) in the same way as we said above in I that \(Y',\allowbreak Y''\), etc., follow from \(Y\); and that each coefficient in \(y\) too can be reduced to the form \[A+B(f,1)+C(f,g)+\text{etc.}\] since the sums of the individual powers of the roots of the equation \(y=0\) are evidently halves of the powers of the roots of the equation \(Y=0\), and hence reducible to such a form. In the last four cases, however, \(Y\) will be the product of the factors \[xx-(\varphi\omega)^2,\ xx-(\varphi \mathfrak{a}\omega)^2,\ xx-(\varphi \mathfrak{b}\omega)^2\text{ etc.}\] and therefore of the form \[x^f-\lambda x^{f-2}+\mu x^{f-4}-\text{etc.}\] It is clear that the coefficients \(\lambda,\allowbreak \mu\), etc., can be derived from the sums of the squares, fourth powers, etc., of the ro\SourceLeaf{465}ots \(\varphi\omega,\allowbreak \varphi \mathfrak{a}\omega,\allowbreak \varphi \mathfrak{b}\omega\), etc.; and the functions \(Y',\allowbreak Y''\), etc., will behave similarly.

\textit{Example} I. Let \(n=17,\allowbreak f=8\) and let \(\varphi\) denote the cosine. We then have \[\begin{aligned}Z=\bigl(&x^8+\tfrac12x^7-\tfrac74x^6-\tfrac34x^5+\tfrac{15}{16}x^4\\&+\tfrac5{16}x^3-\tfrac5{32}xx-\tfrac1{32}x+\tfrac1{256}\bigr)^2\end{aligned}\] and must therefore resolve \(\surd Z\) into two factors of degree four, \(y,\allowbreak y'\). The period \(P=(8,1)\) consists of \((2,1),\allowbreak (2,9),\allowbreak (2,13),\allowbreak (2,15)\), so that \(y\) will be the product of the factors \[x-\varphi\omega,\ x-\varphi9\omega,\ x-\varphi13\omega,\ x-\varphi15\omega\] Substituting \(\tfrac12[k]+\tfrac12[n-k]\) for \(\varphi k\omega\) gives \[\begin{gathered}\varphi\omega+\varphi9\omega+\varphi13\omega+\varphi15\omega=\tfrac12(8,1),\\(\varphi\omega)^2+(\varphi9\omega)^2+(\varphi13\omega)^2+(\varphi15\omega)^2=2+\tfrac14(8,1)\end{gathered}\] Similarly, the sum of the cubes is \(=\tfrac38(8,1)+\tfrac18(8,3)\), and the sum of the fourth powers \(=1\tfrac12+\tfrac5{16}(8,1)\). Determining the coefficients in \(y\) from these by Newton's theorem gives \[\begin{aligned}y&=x^4-\tfrac12(8,1)x^3+\tfrac14((8,1)+2(8,3))xx\\&\quad-\tfrac18((8,1)+3(8,3))x+\tfrac1{16}((8,1)+(8,3))\end{aligned}\] Now \(y'\) is derived from \(y\) by interchanging \((8,1)\) and \((8,3)\); substituting for \((8,1),\allowbreak (8,3)\) the values \(-\tfrac12+\tfrac12\surd17,\allowbreak -\tfrac12-\tfrac12\surd17\) therefore gives \[\begin{aligned}y&=x^4+(\tfrac14-\tfrac14\surd17)x^3-(\tfrac38+\tfrac18\surd17)xx\\&\quad+(\tfrac14+\tfrac18\surd17)x-\tfrac1{16}\end{aligned}\] \[\begin{aligned}y'&=x^4+(\tfrac14+\tfrac14\surd17)x^3-(\tfrac38-\tfrac18\surd17)xx\\&\quad+(\tfrac14-\tfrac18\surd17)x-\tfrac1{16}\end{aligned}\] Similarly, \(\surd Z\) can be resolved into four factors of degree two, of which the first will be \((x-\varphi\omega)(x-\varphi13\omega)\), the second \((x-\varphi9\omega)(x-\varphi15\omega)\), the third \((x-\varphi3\omega)(x-\varphi5\omega)\), and the fourth \((x-\varphi10\omega)(x-\varphi11\omega)\). All the coefficients in these factors can be expressed in terms of the four sums \((4,1),\allowbreak (4,9),\allowbreak (4,3),\allowbreak (4,10)\). Evidently, the product of the first factor and the second will be \(y\), and the product of the third and fourth \(y'\).

\clearpage\textit{Example} II. If, everything else remaining unchanged, \(\varphi\) is supposed to denote the sine, so that \[\begin{aligned}Z&=x^{16}-\tfrac{17}4x^{14}+\tfrac{119}{16}x^{12}-\tfrac{221}{32}x^{10}\\&\quad+\tfrac{935}{256}x^8-\tfrac{561}{512}x^6+\tfrac{357}{2048}x^4-\tfrac{51}{4096}xx+\tfrac{17}{65536}\end{aligned}\] must be resolved into two factors of degree 8, \(y,\allowbreak y'\), then \(y\) will be the product of
\SourceLeaf{466}
four quadratic factors \[xx-(\varphi\omega)^2,\ xx-(\varphi9\omega)^2,\ xx-(\varphi13\omega)^2,\ xx-(\varphi15\omega)^2\] Now, since \(\varphi k\omega=-\tfrac12i[k]+\tfrac12i[n-k]\), we have \[\begin{aligned}(\varphi k\omega)^2&=-\tfrac14[2k]+\tfrac12[n]-\tfrac14[2n-2k]\\&=\tfrac12-\tfrac14[2k]-\tfrac14[2n-2k]\end{aligned}\] From this we derive the sum of the squares of the roots \(\varphi\omega,\allowbreak \varphi9\omega,\allowbreak \varphi13\omega.\varphi15\omega\), namely \(2-\tfrac14(8,1)\), the sum of their fourth powers \(=\tfrac32-\tfrac3{16}(8,1)\), the sum of their sixth powers \(=\tfrac54-\tfrac9{64}(8,1)-\tfrac1{64}(8,3)\), and the sum of their eighth powers \(\tfrac{35}{32}-\tfrac{27}{256}(8,1)-\tfrac1{32}(8,3)\). Hence \[\begin{aligned}y={}&x^8-(2-\tfrac14(8,1))x^6\\&+(\tfrac32-\tfrac5{16}(8,1)+\tfrac18(8,3))x^4\\&-(\tfrac12-\tfrac9{64}(8,1)+\tfrac5{64}(8,3))xx\\&+\tfrac1{16}-\tfrac5{256}(8,1)+\tfrac3{256}(8,3)\end{aligned}\] We derive \(y'\) from \(y\) by interchanging \((8,1),\allowbreak (8,3)\), so that substituting the values of these sums gives \[\begin{aligned}y&=x^8-(\tfrac{17}8-\tfrac18\surd17)x^6+(\tfrac{51}{32}-\tfrac7{32}\surd17)x^4\\&\quad-(\tfrac{17}{32}-\tfrac7{64}\surd17)xx+\tfrac{17}{256}-\tfrac1{64}\surd17\end{aligned}\] \[\begin{aligned}y'&=x^8-(\tfrac{17}8+\tfrac18\surd17)x^6+(\tfrac{51}{32}+\tfrac7{32}\surd17)x^4\\&\quad-(\tfrac{17}{32}+\tfrac7{64}\surd17)xx+\tfrac{17}{256}+\tfrac1{64}\surd17\end{aligned}\] Similarly, \(Z\) can be resolved into four factors whose coefficients can be expressed in terms of sums of four terms; the product of two of them will be \(y\), and the product of the remaining two \(y'\).


